\section{Limitations in full}
\label{app:limitations}

We list these plainly because several of them bound the claims above.

\begin{enumerate}[leftmargin=1.4em,itemsep=2pt,topsep=2pt]
\item \textbf{Learning is not evolution.} Our populations adapt within a training
      run by policy-gradient learning under a fixed reward. The Scott-Phillips
      criterion refers to evolutionary history, and SSI is a within-lifetime
      analogue: it asks whether the emission policy that learning arrives at
      depends on whether receivers respond. A genuine evolutionary outer loop
      over heritable emission traits would test the criterion directly, and we did
      not run one. We write ``emerge under learning'' throughout rather than
      ``evolve''.
\item \textbf{One arena scale, one group size.} Everything is measured with four
      fish in a $60\times\SI{60}{\centi\metre}$ arena, in which the knollenorgan
      channel at unit gain reaches every conspecific. The range sweep partially
      addresses this, but group size is not varied, and the upstream model
      randomises arena size in a way we deliberately do not.
\item \textbf{Pulses are binary and fixed-amplitude.} Mormyrid electric
      communication uses inter-pulse-interval patterning and species- and
      individual-specific waveform structure
      \citep{baker2013multiplexed,carlson2004stereotyped}; in wave-type
      gymnotiforms, amplitude modulations carry honest size information
      \citep{gavassa2012signal}. Our knollenorgan channel is all-or-none, as
      upstream, so waveform and amplitude signalling are structurally impossible
      in the coupled condition. This is precisely the limitation the decoupled
      condition is designed to lift, and it lifts only the crudest version of it:
      one bit of subtype, not a graded waveform.
\item \textbf{Episodes are short.} 512 steps is about \SI{6}{\second} of fish time.
      Dominance relationships in real fish are established over much longer
      horizons, so the dominance-related results should be read as
      within-encounter, not as stable social structure.
\item \textbf{The metabolic cost is a reward term, not a physiological model.} We
      express cost in reward units and report the fraction of the achieved return
      it consumes so that it can be compared with measured energy budgets, but we
      do not model the electrocyte physiology that sets the true cost, nor its
      dependence on amplitude.
\item \textbf{The predator is a caricature.} It is guided by discharges alone. Real
      electroreceptive predators also use vision, mechanoreception and olfaction,
      so real silence buys less safety than it does here, and the cryptic regime is
      correspondingly easier to reach in our model than in a river.
\item \textbf{Regime labels are a classification imposed on a continuum.} The
      phase diagram is generated by thresholding continuous metrics at percentiles
      of their own pooled distribution. The boundaries are therefore descriptive,
      not sharp transitions, and we make no claim about criticality or
      bifurcation.
\end{enumerate}
